\chaptertitle{第三章\quad 两个未知数——从一条线到一个平面}

\sectiontitle{第7讲\quad 方程组——两个条件定两个数}

\subsection*{一个条件不够了}

"鸡和兔关在一个笼子里，一共有 10 个头。"这能算出鸡和兔各几只吗？不能——只知道总头数，有无穷多种可能（1 只鸡 9 只兔，2 只鸡 8 只兔，……）。

但如果再加一个条件："一共有 28 条腿"。鸡 2 条腿，兔 4 条腿。两个条件合在一起，就确定了唯一的答案。

\begin{example}
解鸡兔同笼问题：头共 10 个，腿共 28 条，鸡兔各几只？
\end{example}
\begin{solution}
设鸡 $x$ 只，兔 $y$ 只。
根据头数：$x + y = 10$ \quad (1)\\
根据腿数：$2x + 4y = 28$ \quad (2)

两个方程合在一起，叫\textbf{方程组}。\\

方法一（代入消元）：\\
由 (1) 得 $y = 10 - x$，代入 (2)：\\
$2x + 4(10 - x) = 28$\\
$2x + 40 - 4x = 28$\\
$-2x = -12$，$x = 6$\\
$y = 10 - 6 = 4$\\
答：鸡 6 只，兔 4 只。
\end{solution}

\subsection*{代入消元法}

核心思想：把一个方程变形成"$y$ 等于什么"，代入另一个方程，消去一个未知数。

\begin{example}
解方程组：$\begin{cases}x + y = 10\\2x - y = 5\end{cases}$
\end{example}
\begin{solution}
由 (1) 得 $y = 10 - x$，代入 (2)：\\
$2x - (10 - x) = 5$\\
$2x - 10 + x = 5$\\
$3x = 15$，$x = 5$\\
$y = 10 - 5 = 5$\\
所以 $\begin{cases}x=5\\y=5\end{cases}$。
\end{solution}

\subsection*{加减消元法}

另一种思路：把两个方程相加或相减，消去一个未知数。

\begin{example}
解方程组：$\begin{cases}3x + 2y = 13\\2x + 3y = 12\end{cases}$
\end{example}
\begin{solution}
观察：两个方程中 $x$ 和 $y$ 的系数都不相同，但可以配成相同。

(1)×2：$6x + 4y = 26$\\
(2)×3：$6x + 9y = 36$\\
相减：(6x 抵消) $4y - 9y = 26 - 36$\\
$-5y = -10$，$y = 2$\\
代入 (1)：$3x + 4 = 13$，$x = 3$\\
所以 $\begin{cases}x=3\\y=2\end{cases}$。
\end{solution}

\subsection*{几何意义：两条直线的交点}

$x + y = 10$ 和 $2x - y = 5$ 在坐标平面上各是一条直线。它们的交点 $(5,5)$ 同时满足两个方程——这就是方程组的解。

如果你还没学过直角坐标系也没关系，先记住这个图景：\textbf{方程组是在找两条线的交点}。

\begin{exercise}
\item 用代入消元法解：
  \begin{subexercise}
    \item $\begin{cases}x-y=3\\2x+y=12\end{cases}$
    \item $\begin{cases}x+2y=8\\3x-y=7\end{cases}$
    \item $\begin{cases}3x-2y=7\\x+y=4\end{cases}$
  \end{subexercise}

\item 用加减消元法解：
  \begin{subexercise}
    \item $\begin{cases}2x+3y=12\\3x+2y=13\end{cases}$
    \item $\begin{cases}5x+2y=25\\3x+4y=15\end{cases}$
    \item $\begin{cases}4x-3y=11\\3x+2y=4\end{cases}$
  \end{subexercise}

\item 列方程组解应用：
  \begin{subexercise}
    \item 甲、乙两数的和为 32，甲数的 3 倍等于乙数的 5 倍，求甲、乙两数。
    \item 篮球和排球共 12 个，总价 560 元。篮球每个 50 元，排球每个 40 元，各几个？
    \item 客车和货车共 12 辆，共载 420 人。客车每辆载 40 人，货车每辆载 30 人，各几辆？
  \end{subexercise}

\item 思考：$\begin{cases}x+y=5\\2x+2y=10\end{cases}$ 有多少组解？为什么？
\end{exercise}

\sectiontitle{第8讲\quad 不等式组——两个条件同时成立}

两个不等式放在一起同时满足，就是\textbf{不等式组}。

\begin{example}
解不等式组：$\begin{cases}x+3>0\\2x-1<5\end{cases}$
\end{example}
\begin{solution}
先分别解：
$x+3>0$ → $x>-3$\\
$2x-1<5$ → $2x<6$ → $x<3$\\
两个条件同时成立：$-3 < x < 3$。

在数轴上：$-3$ 和 $3$ 之间，两端都是空心圈。
\end{solution}

\subsection*{不等式组的四种基本类型}

\begin{itemize}
  \item $x>a$ 且 $x>b$（$a>b$）：取 $x>a$（大大取大）
  \item $x<a$ 且 $x<b$（$a<b$）：取 $x<a$（小小取小）
  \item $x>a$ 且 $x<b$（$a<b$）：取 $a<x<b$（大小小大中间找）
  \item $x<a$ 且 $x>b$（$a<b$）：无解（大大小小找不到）
\end{itemize}

\begin{example}
解不等式组：$\begin{cases}2x-1\ge x+1\\x-2\le 2x\end{cases}$
\end{example}
\begin{solution}
$2x-1\ge x+1$ → $x\ge2$\\
$x-2\le 2x$ → $-2\le x$ → $x\ge -2$\\
两个条件同时成立：$x\ge 2$（因为 $2$ 比 $-2$ 大，"大大取大"）。
\end{solution}

\begin{exercise}
\item 解不等式组：
  \begin{subexercise}
    \item $\begin{cases}x+3>0\\2x-1<5\end{cases}$
    \item $\begin{cases}3x-2>4\\x+1<5\end{cases}$
    \item $\begin{cases}2x-1\ge x+1\\x-2\le 2x\end{cases}$
    \item $\begin{cases}x-2<0\\x+1>0\\2x-5\le1\end{cases}$
  \end{subexercise}

\item 含参讨论：$ax<b$ 的解集是什么？（分 $a>0$、$a<0$、$a=0$ 三种情况）

\item 应用：某商品进价 40 元，售价不低于 50 元，且利润不超过进价的 30\%，求售价的取值范围。
\end{exercise}

\sectiontitle{第9讲\quad 一次关系三位一体}

\subsection*{从方程到函数}

到目前为止的路径：\\
$2x+1=7$ → 这是方程，求一个精确的 $x$。\\
$2x+1>7$ → 这是不等式，求 $x$ 的范围。

现在问第三个问题：\textbf{当 $x$ 取不同的值时，$y=2x+1$ 分别等于多少？}

这就是\textbf{函数}——描述输入和输出之间对应关系的规则。

\begin{example}
对于 $y=2x+1$，填写下表并观察规律：
\end{example}
\begin{solution}
\[
\begin{array}{c|cccccc}
x & -3 & -2 & -1 & 0 & 1 & 2 & 3 \\ \hline
y & -5 & -3 & -1 & 1 & 3 & 5 & 7
\end{array}
\]

规律：$x$ 每增加 1，$y$ 增加 2——这就是 \textbf{一次函数}。
\end{solution}

\subsection*{一张图看三位一体}

对于关系 $y = 2x + 1$：

\begin{enumerate}
  \item \textbf{方程} $2x+1=7$：当 $y$ 取定值 $7$ 时，找 $x$。结果是 $x=3$——一个精确值。
  \item \textbf{不等式} $2x+1>7$：当 $y$ 大于 $7$ 时，找 $x$ 的范围。结果是 $x>3$——一个区间。
  \item \textbf{函数} $y=2x+1$：给定任意 $x$，算出对应的 $y$。或者反过来，给定任意 $y$，算出对应的 $x$。
\end{enumerate}

\textbf{三个问法，一个关系。}

\begin{example}
对于 $y = -3x + 6$，求：
\end{example}
\begin{solution}
(1) 当 $y=0$ 时：$-3x+6=0$ → $x=2$（方程，函数图像与 $x$ 轴的交点）

(2) 当 $y>0$ 时：$-3x+6>0$ → $-3x>-6$ → $x<2$（不等式）

(3) 当 $x=4$ 时：$y=-3\times4+6=-6$（函数值）
\end{solution}

\subsection*{一次函数的图像是一条直线}

一次函数 $y=kx+b$ 的图像是一条直线。$k$ 叫\textbf{斜率}，决定直线陡不陡、向上还是向下。$b$ 叫\textbf{截距}，是直线与 $y$ 轴的交点。

\begin{itemize}
  \item 正比例函数 $y=kx$ 是 $b=0$ 的特殊情况，图像过原点。
  \item $k>0$：直线从左下到右上（$x$ 增大 $y$ 增大）
  \item $k<0$：直线从左上到右下（$x$ 增大 $y$ 减小）
\end{itemize}

\begin{example}
指出下列一次函数的斜率和截距：
\end{example}
\begin{solution}
(1) $y=2x+3$：$k=2$，$b=3$（向上倾斜，过 $(0,3)$）

(2) $y=-x+1$：$k=-1$，$b=1$（向下倾斜，过 $(0,1)$）

(3) $y=5x$：$k=5$，$b=0$（正比例，过原点）
\end{solution}

\subsection*{三位一体的几何解释}

对于一条直线 $y=kx+b$：

\begin{itemize}
  \item \textbf{方程} $kx+b=0$：问直线与 $x$ 轴的交点
  \item \textbf{不等式} $kx+b>0$：问直线上方在 $x$ 轴的哪一侧
  \item \textbf{函数} $y=kx+b$：整条直线本身
\end{itemize}

\boxed{\text{方程} \rightarrow \text{交点} \quad \text{不等式} \rightarrow \text{区域} \quad \text{函数} \rightarrow \text{整条线}}

\begin{exercise}
\item 对于 $y=2x-4$，求：
  \begin{subexercise}
    \item 当 $y=0$ 时的 $x$
    \item 当 $y>0$ 时的 $x$ 范围
    \item 当 $y<-2$ 时的 $x$ 范围
    \item 当 $x=3$ 时的 $y$ 值
  \end{subexercise}

\item 对于 $y=-x+3$，求：
  \begin{subexercise}
    \item 当 $y=0$ 时的 $x$
    \item 当 $y\ge0$ 时的 $x$ 范围
    \item 当 $x=-1$ 时的 $y$ 值
    \item 当 $y=x$ 时（即函数值等于自变量）的 $x$
  \end{subexercise}

\item 完成下表并找出规律：
  \begin{subexercise}
    \item $y=3x-2$，填表：$x=-2,-1,0,1,2$ 对应的 $y$
    \item $x$ 每增加 1，$y$ 增加几？
    \item 这个"增加量"和哪个系数有关？
  \end{subexercise}

\item 判断正误并说明理由：
  \begin{subexercise}
    \item 一次函数的图像一定是一条直线。
    \item $y=2x+1$ 与 $y=2x-3$ 的图像平行。
    \item $y=-x+1$ 的图像经过原点。
  \end{subexercise}

\item 思考：$y=2x+1$ 与 $y=2x+1=0$ 有什么区别和联系？
\end{exercise}

\sectiontitle{章末小结}

这一章做了三件事：

\begin{enumerate}
  \item \textbf{方程组}——两个条件确定两个未知数，几何上是两条直线的交点。
  \item \textbf{不等式组}——两个范围条件取公共部分。
  \item \textbf{三位一体}——方程、不等式、函数是同一个关系的三个侧面。
\end{enumerate}

\textbf{核心思想}：一个关系式 $y=kx+b$ 不是方程，不是不等式，也不是函数——它是一个\textbf{关系}。你可以用三种方式问它：
\begin{itemize}
  \item "什么时候等于某值？"→方程
  \item "什么时候大于/小于某值？"→不等式
  \item "每个 $x$ 对应什么 $y$？"→函数
\end{itemize}
